% !TeX root = ../Lua-document-Template.tex
% !TeX spellcheck = uk_UA

\chapter{Постановка задачі повздовжнього зсуву для двошарової прямокутної області}

\section{Математична постановка задачі}

Розглянемо пружне тіло, яке в поперечному перерізі займає прямокутну область $0 < x < a$, $0 < y < b$ і складається з двох шарів (рисунок~\ref{fig:oblast}). Нижній шар $b_0 < y < b_1$ має модуль зсуву $G_1$, верхній шар $b_1 < y < b_2$ — модуль зсуву $G_2$, де $b_0 = 0$, $b_2 = b$. Нижня межа $y = 0$ та бічна межа $x = a$ жорстко закріплені, бічна межа $x = 0$ вільна від напружень, а на верхній межі $y = b$ діє дотичне навантаження $p(x)$.

\begin{figure}[htbp]
	\centering
	\includegraphics[width=0.6\textwidth]{oblast.jpg} % Фото схеми області
	\caption{Двошарова прямокутна область} % Підпис автоматично стане: Рисунок 1.1 — Двошарова прямокутна область
	\label{fig:oblast} % Мітка для посилань у тексті
\end{figure}

В умовах повздовжнього зсуву відмінним від нуля є лише переміщення $w_i(x,y)$ вздовж осі $Oz$, а дотичні напруження в $i$-му шарі визначаються законом Гука:
\begin{equation}
	\tau_{xz}^{(i)} = G_i \frac{\partial w_i}{\partial x}, \qquad
	\tau_{yz}^{(i)} = G_i \frac{\partial w_i}{\partial y}, \qquad i = 1, 2.
	\label{eq:hooke}
\end{equation}

Переміщення в кожному шарі задовольняють рівняння Лапласа:
\begin{equation}
	\frac{\partial^2 w_i}{\partial x^2} + \frac{\partial^2 w_i}{\partial y^2} = 0, \qquad
	0 < x < a, \quad b_{i-1} < y < b_i, \quad i = 1, 2.
	\label{eq:laplace}
\end{equation}

Граничні умови на бічних межах області:
\begin{equation}
	\tau_{xz}^{(i)}\Big|_{x=0} = 0 \;\Longleftrightarrow\; \frac{\partial w_i}{\partial x}\Big|_{x=0} = 0, \qquad
	w_i\Big|_{x=a} = 0, \qquad i = 1, 2.
	\label{eq:side}
\end{equation}

Граничні умови на нижній та верхній межах області:
\begin{equation}
	w_1\Big|_{y=0} = 0, \qquad
	\tau_{yz}^{(2)}\Big|_{y=b} = p(x) \;\Longleftrightarrow\; G_2 \frac{\partial w_2}{\partial y}\Big|_{y=b} = p(x).
	\label{eq:topbottom}
\end{equation}

На межі поділу шарів $y = b_1$ виконуються умови ідеального механічного контакту — неперервність переміщень і дотичних напружень:
\begin{equation}
	w_1\Big|_{y=b_1-0} = w_2\Big|_{y=b_1+0},
	\label{eq:contact_w}
\end{equation}
\begin{equation}
	\tau_{yz}^{(1)}\Big|_{y=b_1-0} = \tau_{yz}^{(2)}\Big|_{y=b_1+0}
	\;\Longleftrightarrow\;
	G_1 \frac{\partial w_1}{\partial y}\Big|_{y=b_1-0} = G_2 \frac{\partial w_2}{\partial y}\Big|_{y=b_1+0}.
	\label{eq:contact_tau}
\end{equation}

Таким чином, задача полягає у знаходженні функцій $w_1(x,y)$ та $w_2(x,y)$, які задовольняють рівняння~(\ref{eq:laplace}), граничні умови~(\ref{eq:side}), (\ref{eq:topbottom}) та умови контакту~(\ref{eq:contact_w}), (\ref{eq:contact_tau}).
